\chapter{The Open Economy}\label{ch:b1:open-economy}

From 1979 to 1998 the Netherlands kept its currency pegged to the Deutsche
Mark, and its short-term interest rate stayed on average $0.46$ points above
Germany's. France and Italy also pegged, in the years the Macrohistory
Database classifies as pegged to the Deutsche Mark (a de facto
classification), and their rates stayed on average $2.76$ and $5.86$ points
above. % figdata: peg_stats
On the afternoon of 16 September 1992 the Swedish central bank raised its
policy rate to 500 per cent to defend its fixed exchange rate. % ledger: b1:riksbank-1992
A country that fixes its exchange rate and lets capital move freely gives
up its own interest rate; a country that tries to keep both may have to pay
such a rate for a few days, and then give up the peg. This chapter opens the
models of the last chapters to trade and capital flows.

\begin{openingchart}{Short-term interest rate minus Germany's (points),
1979--1998, for four countries in the years the Macrohistory Database
classifies as pegged to the Deutsche Mark, a de facto classification (Italy
from 1983, Sweden to 1992).}{Jord\`a--Schularick--Taylor Macrohistory Database, release 6}
\begin{tikzpicture}
\begin{axis}[omtime, xlabel={Year}, ylabel={Gap to Germany (points)}, xmin=1979, xmax=1998, ymin=-2, ymax=14,
  width=0.7\linewidth, height=0.42\linewidth, legend style={font=\scriptsize, at={(0.98,0.98)}, anchor=north east}]
  \addplot[thick, omDef] table[x=year, y=NLD, col sep=comma] {figdata/bachelor-1/27-open-economy/out/pegs.csv};
  \addplot[thick, omThm, dashed] table[x=year, y=FRA, col sep=comma] {figdata/bachelor-1/27-open-economy/out/pegs.csv};
  \addplot[thick, black] table[x=year, y=ITA, col sep=comma, restrict x to domain=1983:1998] {figdata/bachelor-1/27-open-economy/out/pegs.csv};
  \addplot[thick, omRecord, dotted] table[x=year, y=SWE, col sep=comma, restrict x to domain=1979:1992] {figdata/bachelor-1/27-open-economy/out/pegs.csv};
  \legend{Netherlands, France, Italy, Sweden}
\end{axis}
\end{tikzpicture}
\end{openingchart}

The Netherlands' gap was within about a point and a half of zero from 1981
on; France's narrowed only in the 1990s; Italy's stayed between 3 and 12.5
points until 1997, and fell to 1.2 in 1998. % figdata: pegs
In the model of this chapter, a credible peg with free capital flows makes the
home interest rate equal the base country's. Where it did not, the markets
were pricing a devaluation.

\begin{recall}
Exchange rates, the balance of payments, the real exchange rate,
PPP (the hypothesis) and interest parity were met at school. The
gains from trade come from comparative advantage. IS--LM is
\cref{ch:b1:is-lm}; the sector balances are \cref{ch:b1:national-accounts-in-depth}.
\end{recall}

\section{Open-economy accounts}\label{sec:b1:open-economy:accounts}

In an open economy, output is spent at home or abroad:
$Y = C + I + G + NX$, with net exports $NX = X - M$. \omterm{def:b1:long-run-model:saving}{National saving} is $S = Y - C - G$,
so
\[
S - I = NX :
\]
a country that saves more than it invests lends the difference abroad, and
runs a trade surplus. Ignoring income from abroad, $NX$ is the current
account; adding it, the identity becomes $S - I = CA$, with the private,
public and foreign sector balances of \cref{ch:b1:national-accounts-in-depth}.

\begin{definition}[Net foreign asset position]\label{def:b1:open-economy:nfa}
\index{net foreign asset position}
A country's \emph{net foreign asset position} is the value of the assets its
residents hold abroad minus the value of the domestic assets held by
foreigners. A current-account surplus adds to it; so do valuation gains.
\end{definition}

Brenland's accounts for 2025 show net exports of 1,113 million dalens:
\omterm{def:b1:long-run-model:saving}{national saving} of 89,597 million against investment of 88,484 million. Its
government saved $-7{,}865$ million and its private sector 97,462 million. % computed: code/bachelor-1/27-open-economy/exercises/solutions.py

\section{The real exchange rate}\label{sec:b1:open-economy:real}

Let $e$ be the price of a unit of foreign currency in home currency (dalens
per kest), so that a rise in $e$ is a depreciation of the home currency. The
real exchange rate $\varepsilon_r = eP^*/P$ is the price of foreign goods in terms of home
goods; a rise is a real depreciation, which makes home goods cheaper
abroad. Net exports rise with $\varepsilon_r$ --- if the volumes respond enough.

\begin{definition}[Marshall--Lerner condition]\label{def:b1:open-economy:marshall-lerner}
\index{Marshall--Lerner condition}
The \emph{Marshall--Lerner condition} states that, starting from balanced
trade, a real depreciation raises net exports if the price elasticities of
export and import volumes sum to more than one: $\eta_X + \eta_M > 1$.
\end{definition}

In home goods, $NX = X(\varepsilon_r) - \varepsilon_r M(\varepsilon_r)$. Differentiating at $X = \varepsilon_r M$,
\[
\frac{dNX}{d\varepsilon_r} = M(\eta_X + \eta_M - 1),
\]
where $\eta_X = \varepsilon_r X'/X$ and $\eta_M = -\varepsilon_r M'/M$: a depreciation makes each import dearer,
which lowers net exports until the volumes have responded. Since volumes
respond slowly, the condition can fail in the short run and hold later.

\begin{measure}[A real effective exchange rate]
A country trades with many partners. A real effective exchange rate is a
weighted average of its real exchange rates with each of them, the weights
being shares of its trade. The published indices are usually written the
other way round from the convention of this book --- a rise is a real
appreciation --- so a chart must say which.
\end{measure}

\section{The small open economy in the long run}\label{sec:b1:open-economy:long-run}

\begin{definition}[Small open economy]\label{def:b1:open-economy:small-open}
\index{small open economy}
A \emph{small open economy} is one that trades and borrows freely in world
markets and is too small to affect world prices or the world interest
rate: its real interest rate equals the world rate, $r = r^*$.
\end{definition}

With output at its natural level and $r = r^*$, saving and investment no longer
need to be equal; their difference is net exports, $NX = S - I(r^*)$, and the
real exchange rate adjusts so that $NX(\varepsilon_r)$ equals it.

\begin{omfigure}
\begin{tikzpicture}
\begin{axis}[omrecord, grid=none, xlabel={Saving, investment}, ylabel={Real interest rate $r$}, xmin=60, xmax=110,
  ymin=0, ymax=8, width=0.5\linewidth, height=0.4\linewidth, xtick=\empty, ytick=\empty]
  \addplot[supply] coordinates {(89.597,0) (89.597,8)};
  \addplot[demand, domain=63.484:103.484] {(103.484 - x)/5};
  \addplot[guide] coordinates {(60,3) (110,3)};
  \addplot[thick, <->] coordinates {(88.484,3.4) (89.597,3.4)};
  \addplot[shifted] coordinates {(79.597,0) (79.597,8)};
  \node[font=\footnotesize, anchor=south west] at (axis cs:89.8,6.5) {$S$};
  \node[font=\footnotesize, anchor=south east] at (axis cs:79.4,6.5) {$S'$};
  \node[font=\footnotesize, anchor=west] at (axis cs:97,1.2) {$I(r)$};
  \node[font=\footnotesize, anchor=south] at (axis cs:105,3) {$r^*$};
\end{axis}
\end{tikzpicture}
\hspace{1em}
\begin{tikzpicture}
\begin{axis}[omrecord, grid=none, xlabel={Net exports}, ylabel={Real exchange rate $\varepsilon_r$}, xmin=-20, xmax=20,
  ymin=0.6, ymax=1.4, width=0.5\linewidth, height=0.4\linewidth]
  \addplot[demand, domain=-18.887:21.113] {(x + 48.887)/50};
  \addplot[supply] coordinates {(1.113,0.6) (1.113,1.4)};
  \addplot[shifted] coordinates {(-8.887,0.6) (-8.887,1.4)};
  \node[font=\footnotesize, anchor=west] at (axis cs:8,1.18) {$NX(\varepsilon_r)$};
  \node[font=\footnotesize, anchor=south west] at (axis cs:1.3,1.3) {$S - I$};
\end{axis}
\end{tikzpicture}

{\small The \omterm{def:b1:open-economy:small-open}{small open economy} of the weekend problem (billions of dalens).
Left: at the world rate $r^* = 3\%$, investment is 88.484 and saving 89.597; the
difference, 1.113, is lent abroad. Right: the real exchange rate makes net
exports equal to $S - I$. A rise of 10 in government purchases lowers saving
(dashed): net exports fall to $-8.887$ and the real exchange rate appreciates
from 1 to 0.8.}
\end{omfigure}

\begin{model}{The small open economy in the long run}
\index{small open economy}
\assumes $r = r^*$; output at its natural level; saving that does not
depend on $r$.
\predicts A fiscal expansion at home lowers saving, net exports and the
real exchange rate (a real appreciation), and leaves investment
unchanged; a fiscal expansion abroad raises $r^*$, lowers investment at home
and raises net exports.
\tested If capital moved freely, a country's investment would not depend
on its own saving. Across 16 OECD countries in 1960--1974, Feldstein and
Horioka found a coefficient of the investment rate on the saving rate of
$0.89$: most extra saving stayed at home. % ledger: b1:feldstein1980
Across 18 of the 21 countries of their table (they estimated on 16),
averages of 2000--2019 give $0.272$, and
across 49 high-income countries $0.168$ (second chart). % figdata: fh21_stats
% figdata: fh_stats
\breaks Capital is not perfectly mobile; large economies move the world
rate; prices and output adjust in the short run.
\end{model}

\begin{omchart}{Gross investment against gross national saving, per cent of
GDP, average 2000--2019, 49 high-income countries, with the fitted
line.}{World Bank, World Development Indicators}
\begin{tikzpicture}
\begin{axis}[omscatter, xlabel={Gross saving (\% of GDP)}, ylabel={Gross investment (\% of GDP)}, xmin=0, xmax=60,
  ymin=10, ymax=40, width=0.6\linewidth, height=0.42\linewidth]
  \addplot[only marks, mark=*, mark size=1.5pt, omDef] table[x=saving, y=invest, col sep=comma] {figdata/bachelor-1/27-open-economy/out/fh.csv};
  \addplot[omThm, thick, sharp plot, mark=none] table[x=s, y=fit, col sep=comma] {figdata/bachelor-1/27-open-economy/out/fh_fit.csv};
  \addplot[black, thin, dashed, domain=10:40] {x};
  \addplot[only marks, mark=none, nodes near coords, point meta=explicit symbolic, nodes near coords style={font=\scriptsize, anchor=west}]
    table[x=saving, y=invest, meta=iso3, col sep=comma] {figdata/bachelor-1/27-open-economy/out/fh_labels.csv};
\end{axis}
\end{tikzpicture}
\end{omchart}

\begin{record}[Saving and investment]
In 1960--1974, international differences in saving rates among industrial
countries went with almost equal differences in investment rates
(Feldstein and Horioka 1980). % ledger: b1:feldstein1980
With perfect capital mobility the coefficient should have been near zero;
it was $0.89$. Since then the link has weakened: in 2000--2019 it is $0.272$
across 18 of the 21 countries of their table, with a correlation of $0.623$. % figdata: fh21_stats
Norway and Singapore save far more than they invest; Greece and the United
Kingdom invested more than they saved. Capital moves more freely than it
did, but saving still stays largely at home.
\end{record}

\section{The Mundell--Fleming model}\label{sec:b1:open-economy:mf}

In the short run prices are fixed, so $\varepsilon_r$ moves with $e$, and output is set by
demand. With perfect capital mobility, $r = r^*$, and IS--LM becomes:
\[
\text{IS}^*: \quad Y = C(Y - T) + I(r^*) + G + NX(e), \qquad
\text{LM}^*: \quad M/P = L(r^*, Y).
\]
LM$^*$ fixes output; IS$^*$ slopes up in $(Y, e)$: a higher $e$ raises net exports and output.

\begin{model}{The Mundell--Fleming model}
\index{Mundell--Fleming model}
\assumes A \omterm{def:b1:open-economy:small-open}{small open economy}; perfect capital mobility, $r = r^*$; fixed
prices; no expected change in the exchange rate.
\predicts Under a floating rate, a fiscal expansion appreciates the
currency and leaves output unchanged, while a monetary expansion
depreciates it and raises output. Under a fixed rate, the reverse: fiscal
policy moves output by the full multiplier, and monetary policy cannot
move anything.
\tested Interest rates of countries with fixed exchange rates and open
capital markets follow the base country's: over more than 130 years, the
co-movement of interest rates largely supports the constraint
(Obstfeld, Shambaugh and Taylor 2005); the opening chart shows it for
the Netherlands. % ledger: b1:obstfeld2005
\breaks Expected devaluations (Italy's and Sweden's gaps, the rates of
September 1992); prices that adjust; capital that is not perfectly
mobile; large economies.
\end{model}

\begin{proposition}[Policy under the two regimes]\label{prop:b1:open-economy:regimes}
In the Mundell--Fleming model with $C' < 1$ and $NX' > 0$:
\begin{enumerate}
  \item under a float, $dY/dG = 0$, $de/dG = -1/NX'$, and $dY/d(M/P) = 1/L_Y$;
  \item under a peg, $dY/dG = 1/(1 - C')$ and the money supply adjusts by $L_Y\,dY$; a
    change in $M$ has no effect.
\end{enumerate}
\end{proposition}

\begin{proof}
Under a float, LM$^*$ alone gives $Y$, so $G$ cannot move it; IS$^*$ then needs
$NX'de = -dG$. A larger $M/P$ raises $Y$ by $1/L_Y$ and IS$^*$ gives the depreciation. Under a
peg, $e$ is given and IS$^*$ gives $dY = dG/(1 - C')$; to hold $e$ the central bank
supplies the money demanded, $L_Y dY$, buying foreign currency. A purchase of
bonds that raises $M$ would push $r$ below $r^*$; capital would leave, and the bank
would sell foreign currency until $M$ was back where it was.
\end{proof}

\begin{omfigure}
\begin{tikzpicture}
\begin{axis}[omrecord, grid=none, xlabel={Output $Y$}, ylabel={Exchange rate $e$}, xmin=370, xmax=430, ymin=0.6,
  ymax=1.4, width=0.46\linewidth, height=0.4\linewidth, title={\small Float}]
  \addplot[demand, domain=370:430] {(x - 268.265)/125};
  \addplot[shifted, domain=370:430] {(x - 293.265)/125};
  \addplot[supply] coordinates {(393.264,0.6) (393.264,1.4)};
  \addplot[supply, dashed] coordinates {(413.264,0.6) (413.264,1.4)};
  \node[font=\scriptsize, anchor=north] at (axis cs:393.264,1.4) {LM$^*$};
  \node[font=\scriptsize, anchor=west] at (axis cs:412,1.18) {IS$^*$};
\end{axis}
\end{tikzpicture}
\hspace{1em}
\begin{tikzpicture}
\begin{axis}[omrecord, grid=none, xlabel={Output $Y$}, ylabel={Exchange rate $e$}, xmin=370, xmax=430, ymin=0.6,
  ymax=1.4, width=0.46\linewidth, height=0.4\linewidth, title={\small Peg}]
  \addplot[demand, domain=370:430] {(x - 268.265)/125};
  \addplot[shifted, domain=370:430] {(x - 293.265)/125};
  \addplot[black, thick] coordinates {(370,1) (430,1)};
  \addplot[supply] coordinates {(393.264,0.6) (393.264,1.4)};
  \addplot[supply, dashed] coordinates {(418.264,0.6) (418.264,1.4)};
  \node[font=\scriptsize, anchor=south west] at (axis cs:371,1) {peg};
\end{axis}
\end{tikzpicture}

{\small Mundell--Fleming with the weekend problem's numbers, billions of
dalens. IS$^*$ (blue) slopes up: a higher $e$ is a weaker currency. Float: a rise of
10 in $G$ shifts IS$^*$ to the right (dashed) and only lowers $e$ to 0.8; a rise
of 10 in $M/P$ shifts LM$^*$ to 413.264 (dashed red). Peg at $e = 1$: the same rise in $G$
raises output to 418.264, and LM$^*$ follows as the money supply adjusts.}
\end{omfigure}

\section{The trilemma}\label{sec:b1:open-economy:trilemma}

\begin{definition}[Trilemma]\label{def:b1:open-economy:trilemma}
\index{trilemma}\index{impossible trinity}
The \emph{trilemma} (or \emph{impossible trinity}) states that a country cannot
have at once a fixed exchange rate, free capital flows and a monetary
policy of its own: it can have any two.
\end{definition}

Within the model it follows from the condition $r = r^* + \Delta e^e$, where $\Delta e^e$ is the
expected rate of depreciation: with free capital flows any difference in
expected returns is arbitraged away; with a credible peg, $\Delta e^e = 0$, so $r = r^*$.
Monetary policy is then the base country's. Giving up the peg frees $r$;
controls on capital flows break the \omterm{def:b1:price-theory-at-work:arbitrage}{arbitrage}. When the peg is not
credible, $\Delta e^e > 0$ and the home rate must exceed the base rate: the gaps of
France, Italy and Sweden in the opening chart.

\begin{record}[September 1992]
In the summer of 1992 several currencies of the European exchange-rate
mechanism came under attack. On 16 September the Swedish central bank
raised its policy rate to 500 per cent to defend its peg; % ledger: b1:riksbank-1992
the same day, after heavy official purchases of sterling, the United
Kingdom suspended sterling's membership of the mechanism. % ledger: b1:boe-erm-exit-1992
Sweden let its currency float later that autumn. A peg can be held against
an expected devaluation only with interest rates that the rest of the
economy cannot bear for long.
\end{record}

\section*{Exercises}\label{sec:b1:open-economy:exercises}

\begin{exercise}[$\star$]\label{exo:b1:open-economy:nx}
\omterm{def:b1:long-run-model:saving}{National saving} is 400 and investment 350. Find net exports.
\end{exercise}

\begin{exercise}[$\star$]\label{exo:b1:open-economy:real-ex}
The dalen--kest rate is $e = 1.25$ dalens per kest, Kestria's price level 100 and
Brenland's 110. Find the real exchange rate.
\end{exercise}

\begin{exercise}[$\star$]\label{exo:b1:open-economy:read}
From the opening chart's table, give the mean gap to Germany of the
Netherlands, France and Italy in their pegged years. % figdata: peg_stats
Rank them, and say what the ranking suggests about the credibility of
each peg.
\end{exercise}

\begin{exercise}[$\star$]\label{exo:b1:open-economy:predict}
In the \omterm{def:b1:open-economy:small-open}{small open economy} in the long run, predict the effect of a home
fiscal expansion on the real exchange rate.
\end{exercise}

\begin{exercise}[$\star\star$]\label{exo:b1:open-economy:small}
A \omterm{def:b1:open-economy:small-open}{small open economy} has $S = 400$, $I = 500 - 20r$, $r^* = 4$ and $NX = 50\varepsilon_r - 100$. Find
$I$, $NX$ and $\varepsilon_r$.
\end{exercise}

\begin{exercise}[$\star\star$]\label{exo:b1:open-economy:float}
In the Mundell--Fleming model with a floating rate, $G$ rises by 50 and
$NX' = 25$. Find the change in output, the exchange rate and net exports.
\end{exercise}

\begin{exercise}[$\star\star$]\label{exo:b1:open-economy:peg}
Under a peg, the central bank buys bonds to raise $M$. Describe what happens,
step by step, and the final effect.
\end{exercise}

\begin{exercise}[$\star\star$]\label{exo:b1:open-economy:ml}
Starting from balanced trade, the export and import elasticities are 0.6
and 0.7. Does a real depreciation raise net exports? And with 0.3 and 0.4?
\end{exercise}

\begin{exercise}[$\star\star$]\label{exo:b1:open-economy:count}
In the opening chart's table, Italy was pegged for 16 years and Sweden for
14, with mean gaps of $5.86$ and $4.28$ points and largest gaps of $12.53$ and $9.3$.
% figdata: peg_stats
What do these say about the \omterm{def:b1:open-economy:trilemma}{trilemma}, and about the credibility of the pegs?
\end{exercise}

\begin{exercise}[$\star\star\star$]\label{exo:b1:open-economy:derive}
Derive \cref{prop:b1:open-economy:regimes} by differentiating IS$^*$ and LM$^*$.
\end{exercise}

\begin{exercise}[$\star\star\star$]\label{exo:b1:open-economy:sweden}
With an expected devaluation $\Delta e^e$, show that a peg needs $r = r^* + \Delta e^e$. Markets
expect a devaluation of $10\%$ within a week. What annual rate, at simple
interest, does holding the peg require?
\end{exercise}

\begin{exercise}[$\star\star\star$]\label{exo:b1:open-economy:refute}
What evidence would refute the \omterm{def:b1:open-economy:trilemma}{trilemma}?
\end{exercise}

\section{Problem: Brenland Pegs to the Kest}\label{sec:b1:open-economy:problem}

\begin{problem}[{Weekend problem --- accounts, a peg and an attack}]\label{pb:b1:open-economy:peg}
Brenland's accounts for 2025 (invented, millions of dalens) give $Y = 393{,}264$,
$C = 225{,}014$, $I = 88{,}484$, $G = 78{,}653$, $X = 129{,}777$ and $M = 128{,}664$; government revenue was
157,306 and spending 165,171. Exports to Kestria were 40,231 and imports from
it 34,739; the rate was 1.207 dalens per kest in 2024 and 1.242 in 2025.

\medskip\noindent\textbf{Part I --- Accounts.}
\begin{enumerate}
  \item Compute net exports.
  \item Compute \omterm{def:b1:long-run-model:saving}{national saving}.
  \item Check that $S - I = NX$.
  \item Compute \omterm{def:b1:long-run-model:saving}{public saving}.
  \item Compute \omterm{def:b1:long-run-model:saving}{private saving}.
  \item Compute the trade balance with Kestria and with the rest of the world.
  \item By how much did the dalen depreciate against the kest in 2025?
\end{enumerate}

\medskip\noindent\textbf{Part II --- The long run.} In billions: saving is fixed at
89.597, $I = 103.484 - 5r$, $r^* = 3$, and $NX = 50\varepsilon_r - 48.887$.
\begin{enumerate}[resume]
  \item Find $I$, $NX$ and $\varepsilon_r$.
  \item The world rate rises to $4\%$. Find $I$ and $NX$.
  \item Find the new real exchange rate.
  \item Government purchases rise by 10 at $r^* = 3$. Find $NX$ and $\varepsilon_r$.
  \item Why does output not change?
  \item What does the Feldstein--Horioka evidence suggest about $r = r^*$?
\end{enumerate}

\medskip\noindent\textbf{Part III --- Mundell--Fleming.} With the consumption,
investment and money demand of \cref{pb:b1:is-lm:brenland}, $r = r^* = 3$ and
$NX = 50e - 48.887$, with $e$ an index of the exchange rate, 2025 = 1.
\begin{enumerate}[resume]
  \item Write IS$^*$ and LM$^*$, and check the 2025 \omterm{def:b1:models-and-data:equilibrium}{equilibrium}.
  \item Float: $G$ rises by 10. Find $Y$ and $e$.
  \item Float: $M/P$ rises by 10. Find $Y$ and $e$.
  \item Peg at $e = 1$: $G$ rises by 10. Find $Y$.
  \item By how much must the real money supply rise, and how does it happen?
  \item Peg: the central bank raises $M/P$ by 10. What happens?
  \item Which instrument works under which regime?
\end{enumerate}

\medskip\noindent\textbf{Part IV --- An attack.} Markets expect a devaluation of $5\%$
within a month.
\begin{enumerate}[resume]
  \item What must the monthly return on dalen assets exceed?
  \item Annualise the expected depreciation, compounding over twelve months.
  \item Find the annual interest rate needed to hold the peg.
  \item What would that rate do to investment in the toy model?
  \item The central bank holds reserves of 30 billion kests; holders of dalens
    sell 4 billion dalens a day. How many days do the reserves last?
  \item State the result: what rate does holding the peg require?
\end{enumerate}
\end{problem}
